% Paleta del preu (colores del logo)
\definecolor{PreuRosa}{HTML}{E6165D}
\definecolor{PreuGris}{HTML}{A7ABA9}
\definecolor{PreuTinta}{HTML}{3A3D3C}

% Colores de uso en las preguntas.
% En las figuras se usan solo estos nombres, nunca gray!NN ni colores directos,
% para que la misma fuente sirva para las dos ediciones.
%   relleno       superficies claras (fondos de figuras, sectores secundarios)
%   rellenomedio  superficies intermedias (cajas, sectores, regalos, bolitas)
%   rellenofuerte superficies destacadas (barras, sector principal)
%   tinta         trazos oscuros que no son negro puro
\ifedicioncolor
  \colorlet{relleno}{PreuGris!30}
  \colorlet{rellenomedio}{PreuRosa!30}
  \colorlet{rellenofuerte}{PreuRosa!70}
  \colorlet{tinta}{PreuTinta}
  \colorlet{acento}{PreuRosa}
\else
  \colorlet{relleno}{gray!20}
  \colorlet{rellenomedio}{gray!40}
  \colorlet{rellenofuerte}{gray!65}
  \colorlet{tinta}{black!80}
  \colorlet{acento}{black}
\fi

% Texto que cambia según la edición, para enunciados que nombran un color.
% Ejemplo: bolitas \segunedicion{grises}{rosadas} y blancas.
\newcommand{\segunedicion}[2]{\ifedicioncolor#2\else#1\fi}
